\section{Weighted Windows at Split Finite Places}
\label{fw:section}

We extend the geometric criterion by placing locally constant weights
at the selected finite primes. Keep the notation $K/F$, $d=b+2c$,
$\theta=c/d$, $\lambda$ and $\ell=\log\lambda$ of
Section~\ref{geo:section}, where $b>0$ and $K/F$ is unramified at all
finite places. Put $D_0=(\lambda/2)^d$, and keep the archimedean profiles
$f_C,g_D$ and functionals $J_C,J_D$ defined there.

\subsection{The Local Functional}

Let $L$ be a nonarchimedean local field with valuation ring $\mathcal O$,
uniformizer $\pi$, and residue cardinality $Q$. Additive Haar measure gives
$\mathcal O^2$ mass one. Multiplicative Haar measure $du$ on
$\mathcal O^\times$ has mass one, and $n\in\mathbb Z$ has counting measure.
Define
\[
 s(n,u)=(\pi^nu,\pi^{-n}u^{-1}),\qquad
 Tg(z)=\sum_{n\in\mathbb Z}\int_{\mathcal O^\times}g(z+s(n,u))\,du.
\]
For a nonnegative, locally constant, compactly supported $g$, put
\begin{equation}\label{fw:functional}
 A(g)=\int g^p,\qquad Z(g)=\int gTg,\qquad
 \mathcal F_\delta(g)=\frac{Z(g)}{A(g)^{1+\delta}},
 \qquad p=\frac2{1+\delta},\quad0<\delta<1.
\end{equation}
Assume $g\ne0$ and $Z(g)>0$. Only finitely many valuation labels
contribute, since both coordinates of $s(n,u)$ lie in a fixed compact
difference set. The unit average has mass one.

For any integer $a$, put $B_a=\pi^{-a}\mathcal O$. Ball intersection gives
\begin{equation}\label{fw:rectangle-gram}
 \int 1_{B_a\times B_b}\,T1_{B_c\times B_e}
 =Q^{\min(a,c)+\min(b,e)}
   [\max(a,c)+\max(b,e)+1]_+.
\end{equation}
Two translated balls intersect with the smaller volume exactly when
the displacement lies in the larger ball. The first displacement must
have valuation at least $-\max(a,c)$ and the second at least
$-\max(b,e)$. These conditions say
$-\max(a,c)\le n\le\max(b,e)$. Every permitted unit contributes the same
intersection volume. This also proves the formula for negative ball indices.

Put $S_0=\mathcal O$, $S_i=B_i\setminus B_{i-1}$ for $i\ge1$, and
$m_0=1$, $m_i=Q^i-Q^{i-1}$. For $k\ge0$, let $g$ have value $w_{ij}$
on $S_i\times\pi^{-k}S_j$ for $0\le i\le L_x$, $0\le j\le L_y$, and
be zero elsewhere. The weights are nonnegative. With
\[
 D_{ir}=\begin{cases}
 m_{\min(i,r)},&i\ne r,\\
 0,&i=r=0,\\
 i m_i-Q^{i-1},&i=r\ge1,
 \end{cases}
\]
define
\[
 K_{(i,j),(r,t)}=(k+1)m_im_j\mathbf1_{i=r,\,j=t}
       +m_jD_{ir}\mathbf1_{j=t}+m_iD_{jt}\mathbf1_{i=r}.
\]
Then
\begin{equation}\label{fw:shell-kernel}
 A(g)=Q^k\sum_{i,j}m_im_jw_{ij}^p,\qquad Z(g)=Q^k w^TKw.
\end{equation}
To verify the second formula, take the four shell differences of
\eqref{fw:rectangle-gram}. The constant $k+1$ gives the diagonal
term. The double difference of $\max(i,r)Q^{\min(i,r)}$ is $D_{ir}$,
whereas that of $Q^{\min(j,t)}$ is $m_j\mathbf1_{j=t}$. The remaining
term is symmetric. Since $k\ge0$, the positive-part sign can be removed
in all these differences. The formula holds in residue characteristic
two and for ramified $L/\mathbb Q_p$.

These profiles are separately invariant under multiplication by units and
have the additive period
\begin{equation}\label{fw:hard-period}
 C=\mathcal O\times\pi^{-k}\mathcal O,\qquad |C|=Q^k,
\end{equation}
irrespective of the number of shells. In particular $w_{00}=1$ and all
other weights zero recover $\mathcal F_\delta=(k+1)Q^{-\delta k}$.

For product weights $w_{ij}=a_i b_j$, define
\[
 P_a=\sum_i m_i a_i^p,\quad S_a=\sum_i m_i a_i^2,\quad
 R_a=\sum_{i,r}a_iD_{ir}a_r,
\]
and define $P_b,S_b,R_b$ similarly. Formula
\eqref{fw:shell-kernel} becomes
\begin{equation}\label{fw:product-shells}
 A(g)=Q^kP_aP_b,\qquad
 Z(g)=Q^k[(k+1)S_aS_b+R_aS_b+S_aR_b].
\end{equation}
Thus symmetric product weights need only one vector. The one-shell choice
$(a_0,a_1)=(b_0,b_1)=(1,x)$ has
$S=1+(Q-1)x^2$ and $R=2x+(Q-2)x^2$, so its log gain over the hard
rectangle is
\begin{equation}\label{fw:one-shell-gain}
 \log\frac{(k+1)S^2+2SR}{k+1}
 -2(1+\delta)\log(1+(Q-1)x^p).
\end{equation}
Its right derivative at zero is $4/(k+1)>0$. For the numerical
certificate we use the full finite functional.

\subsection{Finite Periods and Lattice Counting}

Let $S$ now denote a finite set of $F$-primes split in $K$. At each
$v\in S$ identify the two completions with $L_v=F_v$, compatibly with
$\iota$, and choose a common uniformizer. Put
\[
 E_S=\prod_{v\in S}(L_v\times L_v),\qquad
 \Gamma=\mathcal O_{K,S}\subset\mathbb C^d\times E_S.
\]
Here $\mathcal O_{K,S}$ consists of the elements integral outside the
selected prime pairs. For a rectangular compact open subgroup
$C_f=\prod_v(\pi_v^{a_v}\mathcal O_v\times\pi_v^{b_v}\mathcal O_v)$,
the elements of $\Gamma$ with finite component in $C_f$ form the fractional
ideal
\[
 L=\prod_v\mathfrak P_v^{a_v}(\iota\mathfrak P_v)^{b_v},\qquad
 N_K(L)=|C_f|^{-1},\qquad \operatorname{covol}(L)=D_0/|C_f|.
\]
Every $C_f$-coset in $E_S$ contains a finite component of a $\Gamma$-point.
To see this, enclose the coset and $C_f$ in larger local balls and form
the corresponding fractional ideal $J\supset L$. After clearing a common
denominator, the Chinese Remainder Theorem identifies $J/L$ with its
local quotients at the selected prime ideals. This requires no
principality hypothesis, and each finite coset has an archimedean fiber
equal to a translate of $L$.

If $P_L$ is an archimedean fundamental parallelepiped, then $P_L\times C_f$
is a fundamental domain for $\Gamma$, of Haar volume $D_0$. First subtract
a $\Gamma$-element to put the finite component in $C_f$, then subtract an
$L$-element to put the infinite component in $P_L$. The same argument
proves uniqueness and discreteness. Reciprocal archimedean deformation
$B_h$ gives $\Gamma_h=(B_h,\mathrm{id})\Gamma$ with the same covolume.

Write $H_C=d^{-1}\log|C_f|$. The codifferent argument in
\eqref{geo:dual-product} applied to $L$ gives the product minimum
$[2e^{H_C/2-\ell}]^d$ on every deformed dual lattice. Consequently the
archimedean Fourier condition sufficient for this section is
\begin{equation}\label{fw:fourier-condition}
 \sigma\,2e^{H_C/2-\ell}\ge\log M+2\log5+2,
\end{equation}
using the envelope and pointwise regularity of Section~\ref{geo:section}.
It gives a bound $2\int F_\infty/(D_0/|C_f|)$ for the sum of
$F_\infty=\prod f_C^p\prod g_D^p$ on every archimedean fiber.
Summing against any nonnegative, compactly supported, $C_f$-periodic
function $F_f$ yields
\begin{equation}\label{fw:weighted-count}
 \sum_{\gamma\in\Gamma_h+w}F_\infty(\gamma_\infty)F_f(\gamma_f)
 \le\frac2{D_0}\int F_\infty\int F_f.
\end{equation}
The values of $F_f$ on its finitely many nonzero cells sum to
$(\int F_f)/|C_f|$, so the estimate depends on the period volume.
Local reciprocal diagonal rescaling replaces $C_f$ by a rectangular
period of the same volume. Hence the bound remains uniform after such
rescalings, including those that do not preserve $\Gamma$.

\subsection{Class Selection and Unit Averaging}

Let $\Omega\subset\mathbb C^d\times E_S$ be bounded, have compact finite
support, and be invariant under all coordinate rotations and local
reciprocal unit multiplications. Denote its overlap with a norm-one step
of archimedean pair-log coordinates $u\in\mathbb R^c$ and finite valuation
labels $n\in\mathbb Z^S$ by $\Phi(u,n)$. Invariance removes all phase and
unit choices. Only finitely many labels occur, and the $u$-support is
bounded. Form the weights
\[
 W(n)=\int_{\mathbb R^c}\Phi(u,n)\,du.
\]
The quotient $\mathrm{Cl}(K)/\mathrm{im}\,\mathrm{Cl}(F)$ has order
$\kappa h_{\rm rel}$. Map $n$ to the class of
$A(n)=\prod_v\mathfrak P_v^{n_v}$. One class has total weight at least
$\sum_nW(n)/(\kappa h_{\rm rel})$. Suppose this sum is positive and
choose a positive-weight reference label $n_0$ in that class.
For every selected label, write
$A(n)/A(n_0)=(x_n)\mathfrak a_n\mathcal O_K$. Then
\[
 \beta_n=x_n/\iota(x_n),\quad N_{K/F}\beta_n=1,\quad
 (\beta_n)=\prod_v\mathfrak P_v^{n_v-n_{0,v}}
                  (\iota\mathfrak P_v)^{-n_v+n_{0,v}}.
\]
Then $\beta_n\in\mathcal O_{K,S}$ and the full set of norm-one generators
of this ideal is $\beta_nU$. Distinct labels give disjoint cosets.

Let $\alpha_0$ be the identity at infinity with finite coordinates
$(\pi_v^{n_{0,v}},\pi_v^{-n_{0,v}})$, and put $\Omega'=\alpha_0^{-1}\Omega$.
This preserves Haar and period volumes. The overlap with the displacement
$(B_h\beta,\beta_f)$, for $\beta\in\beta_nU$, is
$\Phi(l(\beta)-h,n)$. Since the overlap is independent of units and
phases, correlations between local units of global generators and their
archimedean data do not affect it. Tiling by the full unit lattice gives
\[
 \frac1{R_U}\int_{P_U}\sum_{n\text{ selected}}
 \sum_{\beta\in\beta_nU}\Phi(l(\beta)-h,n)\,dh
 =\frac{w_K}{R_U}\sum_{n\text{ selected}}W(n).
\]
For each fixed field and bounded $h$-domain, only finitely many unit-log
lattice points meet the bounded overlap support, and the logarithmic
kernel has $w_K$ elements. Tonelli also justifies the sums directly,
since their terms are nonnegative. We need no distribution hypothesis
on the local unit images.

Jointly average $h$ and additive translations over a fundamental domain
of $\Gamma_h$, then weight parameter space by its point count. The
translation averages are $|\Omega|/D_0$ for vertices and the sum of
overlaps divided by $D_0$ for ordered edges. Hence some pair of
transformations gives a nonempty point set satisfying
\begin{equation}\label{fw:joint-average}
 \frac EN\ge
 \frac{w_K}{\kappa h_{\rm rel}R_U}
 \frac{\sum_n\int\Phi(u,n)\,du}{|\Omega|}.
\end{equation}
The uniform bound \eqref{fw:weighted-count} holds for this same pair.
Choose the parameter domains using a continuously deformed basis and
Haar measure on the finite period. Countable indicator sums and Tonelli
justify the averages, including sublevel boundaries, since the
exponential majorant is pointwise and additive tiling holds for
measurable indicators. The class is chosen from the overlaps of the
final window, so concentration will be applied before class selection.

\subsection{The Finite-Place Transfer}

For every selected place choose a nonzero, nonnegative, locally constant,
compactly supported $g_v$ with a rectangular additive period $C_v$ and
the unit invariance above, and with $Z_v=Z(g_v)>0$. Suppose their number
is $O(d)$ and their profiles belong to a fixed finite list of local types.
The finite shell profiles above meet these hypotheses. Write
$A_v=A(g_v)$ and $\mathcal F_{\delta,v}=Z_v/A_v^{1+\delta}$. These local
functionals give the following extension of Proposition~\ref{geo:transfer}.

\begin{proposition}\label{fw:transfer}
Suppose the field degrees tend to infinity, $\log L_F(1)\le dC$, and
\eqref{fw:fourier-condition} holds with $C_f=\prod_v C_v$.
Assume the archimedean energy and moment hypotheses of
Proposition~\ref{geo:transfer}. If
\begin{equation}\label{fw:margin}
 \begin{split}
 \mathcal M_{\rm fin}(\theta)={}&
 \frac1d\sum_{v\in S}\log\mathcal F_{\delta,v}
 -(1/2-\delta)\ell-C+(1-\delta)\log2+(1-\theta)\log\pi\\
 &+(1-2\theta)J_C+\theta J_D
 \end{split}
\end{equation}
has a fixed positive lower bound, there are planar sets of cardinality
tending to infinity whose unordered unit-distance count divided by
$n^{1+\delta}$ tends to infinity.
\end{proposition}

\begin{proof}
The total profile $f=\prod f_C\prod g_D\prod_v g_v$ has
\[
 A=\int f^p=A_C^bA_D^c\prod_v A_v,\qquad
 Z=I_C^bI_D^c\prod_v Z_v.
\]
Use the independent normalized overlap probability laws of all these
blocks. The two endpoint energy distributions agree by reflection and
the stated symmetries. Let $d\bar v$ be the sum of first-endpoint means
for energy $-\log f$. On each finite factor, the positive support is a
finite union of additive cells with a positive minimum weight, so its
endpoint energy is bounded. The total variance is $O(d)$.

On the positive support put
$\Omega=\{-\log f\le d(\bar v+\varepsilon)\}$, for $\varepsilon>0$.
Coercivity makes it bounded. Its finite support is compact and it has
all the required periods and invariances. Exponential majorization and
\eqref{fw:weighted-count}, also after the reference-label rescaling,
give
\[
 |\Omega|\le A e^{pd(\bar v+\varepsilon)},\qquad
 N\le 2A e^{pd(\bar v+\varepsilon)}/D_0.
\]
Chebyshev's inequality and a union bound put both endpoint energies in
$[d(\bar v-\varepsilon),d(\bar v+\varepsilon)]$ with probability at least
$1/2$ for large $d$. Undoing the overlap density on this event gives
\[
 \sum_n\int\Phi(u,n)\,du\ge\tfrac12 Z e^{2d(\bar v-\varepsilon)}.
\]
Since $p(1+\delta)=2$, \eqref{fw:joint-average} implies
\[
 \frac E{N^{1+\delta}}\ge
 \frac{w_K D_0^\delta}{2^{1+\delta}\kappa h_{\rm rel}R_U}
 \frac{Z}{A^{1+\delta}}e^{-4\varepsilon d}.
\]
Substituting \eqref{geo:mass-density}, the ordered ratio is at least
$2^{-2-\delta}e^{d(\mathcal M_{\rm fin}-4\varepsilon)}$.
Choose $4\varepsilon<\inf\mathcal M_{\rm fin}$ and divide by two for
unordered edges. Projection to a compact coordinate is injective on the
field points and sends the norm-one displacements to unit vectors.
The upper bound $N^2/2$ for unordered edges forces $N\to\infty$.
\end{proof}

If a rational prime has actual indices $e_p,f_p$ and all its $F$-places
split in $K/F$, there are $d/(e_pf_p)$ such places. A common local profile
contributes $\log\mathcal F_{\delta,p}/(e_pf_p)$ to
\eqref{fw:margin}. Hard rectangles give
$\log\mathcal F_{\delta,p}=\log(k_p+1)-\delta k_pf_p\log p$, so their
sum is exactly $J-\delta H$. Adding finite shells with period
\eqref{fw:hard-period} preserves
$H_C=\sum_p k_p\log p/e_p=H$. Hence the finite-place functional replaces
the ideal-choice term and leaves the Fourier condition unchanged.
